%% kapa: exam=1 round=26 year=2015 period=2 form=A subject=appraisal_law subject_ko=부동산관계법규 seq=24 number=24
%% source: Q-Net 감정평가사 기출문제 (cst00309) artlSeq=5138001 file=extracted/1008/5138001/제26회_감정평가사 1차_기출/제26회_감정평가사 1차_기출/2교시_감평_A형.pdf
%% conversion: kordoc 4.24.51 -> tools/gen1.py (auto); review=proofread-sonnet; flags=none
\begin{question}{24}
부동산가격공시 및 감정평가에 관한 법령상 개별공시지가에 관한 설명으로 옳지 않은 것은?
\begin{choices}
  \item 시장ㆍ군수 또는 구청장은 개별공시지가에 토지가격비준표의 적용에 오류가 있는 경우 시ㆍ군ㆍ구 부동산평가위원회의 심의를 거치지 아니하고 직권으로 정정하여 결정ㆍ공시하여야 한다.
  \item 개별공시지가를 공시하는 시장ㆍ군수 또는 구청장은 필요하다고 인정하는 때에는 개별공시지가의 결정 및 이의신청에 관한 사항을 토지소유자등에게 개별통지할 수 있다.
  \item 국토교통부장관은 개별공시지가의 조사ㆍ산정지침을 정하여 시장ㆍ군수 또는 구청장에게 통보하여야 한다.
  \item 시장ㆍ군수 또는 구청장이 개별공시지가를 결정ㆍ공시하는 경우에는 당해 토지와 유사한 이용가치를 지닌다고 인정되는 하나 또는 둘 이상의 표준지의 공시지가를 기준으로 토지가격비준표를 사용하여 지가를 산정하여야 한다.
  \item 국토교통부장관은 지가공시행정의 합리적인 발전을 도모하고 표준지공시지가와 개별공시지가와의 균형유지 등 적정한 지가형성을 위하여 필요하다고 인정하는 경우에는 개별공시지가의 결정ㆍ공시 등에 관하여 시장ㆍ군수 또는 구청장을 지도ㆍ감독할 수 있다.
\end{choices}
\end{question}
\answer{1}
